\chapter{Nearest Neighbors and Naive Bayes}
\label{ch:knn-nb}

\chapterguide%
  {Distance measures from \chapref{ch:linear-algebra}, Bayes' theorem from \chapref{ch:statistics}, and model-selection ideas from \chapref{ch:training}.}
  {apply nearest neighbors and naive Bayes; explain scaling, smoothing, and their assumptions.}

$k$-nearest neighbors predicts from similar examples; naive Bayes compares how well each class explains the observed features. Both are useful baselines with hand-checkable calculations.

\section{$k$-nearest neighbors}
\label{sec:knn}

$k$-NN has no training phase worth the name: it simply \emph{stores} the training data. All the work happens at prediction time:

\begin{enumerate}
  \item compute the distance from the new observation to every stored training observation;
  \item keep the $k$ closest ones --- the \term{neighbors};
  \item for classification, predict the majority class among the neighbors; for regression, predict the average of their labels.
\end{enumerate}

\begin{mlfigure}
\begin{tikzpicture}[scale=1.05, font=\small]
\foreach \p in {(1,1),(1.5,2),(2,1.2),(0.8,2.2)}
  \fill \p circle (2.3pt);
\foreach \p in {(3.5,3),(4,2.5),(3,3.6),(4.3,3.4)}
  \draw[thick, fill=gray!35] \p circle (2.3pt);
\node[star, star points=5, star point ratio=2.4, fill=black,
      inner sep=1.6pt] (q) at (2.6,2.4) {};
\draw[dashed] (2.6,2.4) circle (1.32);
\node[right=2pt] at (2.78,2.4) {new point};
\node[align=left] at (1.1,3.5) {$\bullet$ = class A};
\node[align=left] at (4.6,1.5) {$\circ$ = class B};
\node[align=center] at (2.6,0.55)
  {$k=3$ neighborhood: two of the three\\nearest points are class B $\rightarrow$ predict B};
\end{tikzpicture}
\end{mlfigure}

\begin{workedexample}
Training data: $A=(1,1)$ and $B=(2,1)$ labeled Red; $C=(4,4)$ and $D=(5,4)$ labeled Blue. Classify the new point $(2,2)$ with $k=3$ using Euclidean distance:
\begin{gather*}
  d_A=\sqrt{1^2+1^2}\approx 1.41,\qquad
  d_B=\sqrt{0^2+1^2}=1.00,\\
  d_C=\sqrt{2^2+2^2}\approx 2.83,\qquad
  d_D=\sqrt{3^2+2^2}\approx 3.61.
\end{gather*}
The three nearest are $B$ (Red), $A$ (Red), and $C$ (Blue). The vote is Red 2, Blue 1, so the prediction is Red.
\end{workedexample}

\subsection{Choosing $k$}

The single hyperparameter $k$ is a smoothness dial:

\begin{itemize}
  \item \term{Small $k$} gives a flexible boundary that follows local patterns but also noise: low bias, high variance.
  \item \term{Large $k$} smooths the boundary but can blur distinctions. At $k=n$, classification always predicts the majority class: high bias, low variance.
\end{itemize}

Validate $k$. Odd $k$ avoids binary majority-vote ties; distance weighting gives nearer neighbors more influence.

\subsection{What $k$-NN is sensitive to}

\begin{mltable}
\begin{tabularx}{\linewidth}{@{}>{\bfseries\leavevmode\color{MLNavy}}p{0.20\linewidth}p{0.42\linewidth}X@{}}
\toprule
\rowcolor{MLPaleBlue!92}
Sensitivity & Why it matters & Practical response \\
\midrule
Feature scale & Distance inherits the unit problem from \secref{sec:scaling}: an unscaled \texttt{area} column would decide every neighbour on its own. & Standardize inside the pipeline, as Labs~1 and~3 do. \\
Irrelevant features & Every feature affects the distance, so noise features dilute the signal. & Feature selection helps $k$-NN more than most models. \\
High dimensions & With many weak or noisy features, distance contrast can shrink --- the \term{curse of dimensionality} --- and ``nearest'' may stop meaning ``similar.'' & Feature selection or dimensionality reduction (\chapref{ch:dimred}) can help. \\
Dataset size & Predicting requires comparing against the whole training set. & Tree-based indexes (KD-trees, ball trees) and approximate nearest-neighbor methods make this workable at scale. \\
\bottomrule
\end{tabularx}
\end{mltable}

$k$-NN also does regression --- predict the average (or distance-weighted average) label of the neighbors --- which yields a locally adaptive, assumption-free fit, at the price of the same sensitivities listed above.

\begin{intuition}
$k$-NN never builds a theory of the data; it just answers every question with ``let me check what similar past cases did.'' That makes it honest and flexible --- and completely dependent on the distance actually measuring similarity that matters.
\end{intuition}

\section{How $k$ reshapes the boundary}

The boundary changes as $k$ changes.

\begin{mlfigure}
\begin{tikzpicture}[scale=0.88, every node/.style={font=\scriptsize}]
% k=1
\begin{scope}
\draw[draw=MLMuted,fill=white] (0,0) rectangle (3.4,2.7);
\draw[MLRed,very thick,rounded corners=1pt]
  (0,1.15) -- (0.5,1.4) -- (0.75,0.95) -- (1.2,1.5) -- (1.45,1.0) -- (1.9,1.6)
  -- (2.15,1.05) -- (2.6,1.65) -- (2.9,1.1) -- (3.4,1.5);
\foreach \p in {(0.4,1.9),(0.9,2.2),(1.5,2.0),(2.2,2.3),(2.8,1.95),(1.9,2.45)}
  \fill[MLBlue] \p circle (2.1pt);
\foreach \p in {(0.5,0.5),(1.1,0.4),(1.7,0.6),(2.4,0.35),(2.9,0.6),(1.4,0.15)}
  \fill[MLOrange] \p circle (2.1pt);
\fill[MLOrange] (1.15,1.62) circle (2.1pt);
\node[font=\small\bfseries\color{MLNavy}] at (1.7,3.05) {$k=1$};
\node[color=MLMuted] at (1.7,-0.38) {jagged, chases every point};
\end{scope}
% k=15
\begin{scope}[xshift=4.4cm]
\draw[draw=MLMuted,fill=white] (0,0) rectangle (3.4,2.7);
\draw[MLGreen,very thick] (0,1.2) .. controls (1.2,1.45) and (2.2,1.15) .. (3.4,1.4);
\foreach \p in {(0.4,1.9),(0.9,2.2),(1.5,2.0),(2.2,2.3),(2.8,1.95),(1.9,2.45)}
  \fill[MLBlue] \p circle (2.1pt);
\foreach \p in {(0.5,0.5),(1.1,0.4),(1.7,0.6),(2.4,0.35),(2.9,0.6),(1.4,0.15)}
  \fill[MLOrange] \p circle (2.1pt);
\fill[MLOrange] (1.15,1.62) circle (2.1pt);
\node[font=\small\bfseries\color{MLNavy}] at (1.7,3.05) {$k=15$};
\node[color=MLMuted] at (1.7,-0.38) {smooth, ignores the odd point};
\end{scope}
\end{tikzpicture}
\end{mlfigure}

\begin{intuition}
At $k=1$, the boundary forms an island around the isolated orange point. At $k=15$, surrounding blue votes outweigh it. Smoothing reduces sensitivity to isolated noise.
\end{intuition}

\subsection{SMOTE: nearest neighbours used to invent minority examples}
\label{sec:smote}

\chapref{ch:training} listed SMOTE as one way to handle a rare class. Its mechanics are pure $k$-NN. Take one minority-class row, find its $k$ nearest \emph{minority} neighbours, pick one of them at random, and create a new synthetic row somewhere on the straight line between the two. Repeat until the classes are as balanced as you want.

SMOTE inherits distance-scaling sensitivity. Interpolation can create invalid categorical or binary values, so use appropriate representations and methods:

\begin{itemize}
  \item scale the features first, inside the pipeline;
  \item split first, and synthesize only from the training portion of each fold, never from held-out rows.
\end{itemize}

\section{Naive Bayes}
\label{sec:naive-bayes}

Naive Bayes turns classification into the question Bayes' theorem was built for: given the observed features, which class is most probable?
\[
  P(y\given x_1,\ldots,x_p)
  \;\propto\;
  P(y)\prod_{j=1}^{p}P(x_j\given y).
\]
$P(y)$ is the class \term{prior}; $P(x_j\given y)$ is a feature \term{likelihood}. Multiplying likelihoods assumes conditional independence (\chapref{ch:statistics}). Though often unrealistic, this can still rank classes usefully.

Predict $\hat y=\argmax_y P(y)\prod_jP(x_j\given y)$. Compute in log space to avoid underflow.

\subsection{Three standard variants}

\begin{mltable}
\begin{tabularx}{\linewidth}{@{}>{\bfseries\leavevmode\color{MLNavy}}p{0.20\linewidth}p{0.25\linewidth}X@{}}
\toprule
\rowcolor{MLPaleBlue!92}
Variant & Feature type & How it models the features \\
\midrule
Gaussian NB & Continuous features & Each $P(x_j\given y)$ is a normal density with the class-specific mean and variance of feature $j$. \\
Multinomial NB & Count data, such as word counts & $P(x_j\given y)$ is the probability of word $j$ in class $y$'s text, and repeated words count repeatedly. It is the workhorse for text classification. \\
Bernoulli NB & Binary presence/absence features & Explicitly models both a word appearing and a word being absent. \\
\bottomrule
\end{tabularx}
\end{mltable}

\subsection{Laplace smoothing}

Raw counts assign zero probability to unseen class--word pairs, eliminating the whole product. \term{Laplace (add-one) smoothing} adds one to every count:
\[
  P(w\given y)
  =\frac{\operatorname{count}(w,y)+1}
        {\operatorname{total\ words\ in\ }y+V},
\]
where $V$ is vocabulary size. Smoothing prevents zero probabilities.

\begin{workedexample}
A tiny spam filter with vocabulary $\{$money, free, meeting, lunch$\}$ ($V=4$) is trained on 10 words of spam and 10 words of ham, with equal priors $P(\text{spam})=P(\text{ham})=0.5$:

\begin{mltable}
\begin{tabular}{@{}lcccc@{}}
\toprule
\rowcolor{MLPaleBlue!92}
raw counts & money & free & meeting & lunch \\
\midrule
spam (10 words) & 5 & 4 & 1 & 0 \\
ham (10 words) & 1 & 1 & 5 & 3 \\
\bottomrule
\end{tabular}
\end{mltable}

With smoothing, every likelihood is $(\text{count}+1)/(10+4)$; for example $P(\text{money}\given\text{spam})=6/14$ and $P(\text{money}\given\text{ham})=2/14$. Classify the message ``free money'':
\[
  \text{spam score}=0.5\times\tfrac{6}{14}\times\tfrac{5}{14}
  =\tfrac{15}{196},
  \qquad
  \text{ham score}=0.5\times\tfrac{2}{14}\times\tfrac{2}{14}
  =\tfrac{2}{196}.
\]
Normalizing gives spam probability $15/(15+2)\approx0.88$. Without smoothing, any message containing ``lunch'' would receive zero spam probability.
\end{workedexample}

\subsection{Strengths and limitations}

\begin{itemize}
  \item Per-feature statistics make training and prediction efficient, especially for sparse text.
  \item Few estimated interactions make it effective with limited data.
  \item Correlated evidence can produce overconfident probabilities; validate calibration when risk estimates matter (\chapref{ch:evaluation}).
  \item Conditional independence prevents explicit modeling of feature interactions.
\end{itemize}

\begin{keyidea}
$k$-NN and naive Bayes make opposite simplifications. $k$-NN keeps the training examples rather than estimating a fixed parametric form, but it still assumes that the chosen distance and local neighborhood are meaningful. Naive Bayes compresses the data into per-class, per-feature summaries and assumes conditional independence. Both are quick to build, valuable as baselines, and clear about what they ignore.
\end{keyidea}

\begin{modelcard}{$k$-NN and Naive Bayes - Quick Guide}
\begin{tabularx}{\linewidth}{@{}>{\bfseries\leavevmode\color{MLNavy}}p{0.22\linewidth}X X@{}}
\toprule
\rowcolor{MLPaleBlue!92}
 & $k$-nearest neighbors & Naive Bayes \cr
\midrule
Core idea & Predict from nearby training examples & Combine class prior with per-feature evidence \cr
Best when & Local similarity is meaningful and data is not too large & Fast probabilistic classification is needed, especially for count or text features \cr
Main controls & Number of neighbors, distance, weighting & Distribution variant, smoothing strength \cr
Main risks & Scale, irrelevant features, slow prediction, high dimensions & Unrealistic independence assumption, probability calibration \cr
\bottomrule
\end{tabularx}
\end{modelcard}

\begin{quickcheck}
\begin{enumerate}
  \item Why is feature scaling usually essential for KNN, and what role does \(k\) play in the bias--variance trade-off?
  \item Why can KNN become expensive at prediction time even though its training step is simple?
\end{enumerate}
\end{quickcheck}
